\cotefile{156-7}
\scanpage{73}
\mspage{72}

on a une notion de ``support \unsure{minimal}'' \unsure{subordonné} à cette figure, \unsure{l'ens.|l'ens} \inserted{\unsure{ordonné}} de ses supports s'identifie à $\unsure{\mathcal{P}|\mathscr{P}}(\widetilde F)$. Et si $F'\unsure{\preceq|\leq}F$ est une subdivision, plus gén. $F'\ll F$ un raffinement, on \unsure{sent|sait} comment ``passer'' de l'un à l'autre, via
\[
\unsure{\varphi|\phi}:\widetilde{F'}\longrightarrow\widetilde F
\]
induisant
\[
\varphi^*:\mathcal{P}(\widetilde F)\longrightarrow\mathcal{P}(\widetilde{F'}).
\]

\unsure{Il} ne \unsure{semblerait} pas, dans le cas d'un atlas qui ne soit \unsure{quasi-constructible}, qu'on puisse interpréter \deleted{\unsure{donc}} de façon canonique les $\mathcal{P}(\widetilde F)$, pour des figures \unsure{différentes|differentes} (et qui ne seraient ni compatibles, ni \unsure{reliées|reliés} par $F'\ll F$, ni \inserted{pré-}compatibles dans le sens qu'elles deviennent compatibles après \unsure{subd.|subdiv.} de l'une et de l'autre\ldots), comme des parties d'un même ensemble $\mathcal{P}(P)$ disons, où $P$ est un ensemble de ``points'' associé d'une façon ou d'une autre au \deleted{\unsure{M.}} Atlas $M$.

On va définir cependant la notion de ``réalisation par supports'' d'un \unsure{stratos|Atlas}, dans un ensemble $P$ de ``points''.
